\documentclass{thapathali}

% ============================================
% CUSTOMIZABLE VARIABLES - Lean Software Development
% ============================================
\projecttitle{LEAN SOFTWARE DEVELOPMENT}
\projectsubtitle{Core Philosophy, Principles, and Implementation}
\newcommand{\studentone}{Prasiddha Raj Poudel}
\newcommand{\rollone}{THA080BCT027}
\newcommand{\studenttwo}{Pratik Pant}
\newcommand{\rolltwo}{THA080BCT028}
\newcommand{\studentthree}{Purnima Bhandari}
\newcommand{\rollthree}{THA080BCT029}
\newcommand{\studentfour}{}
\newcommand{\rollfour}{}

% Used in structural settings
\presentername{\studentone}
\studentroll{\rollone}
\supervisorname{Dr. Supervisor Name}
\supervisortitle{Associate Professor}
\departmentname{Department of Electronics and Computer Engineering}
\institutionname{Institute of Engineering, Thapathali Campus}
\presentationdate{July 9, 2026}
\academicyear{Academic Year: 2026/2027}

% ============================================
% TITLE PAGE SETUP
% ============================================
\csname title\endcsname[\ProjectTitle]{\ProjectTitle}
\subtitle{\ProjectSubtitle}
\author[\studentone]{\studentone, \studenttwo, \studentthree}
\institute[\DepartmentName]{\DepartmentName}
\date[\PresentationDate]{\PresentationDate}

% ============================================
% DOCUMENT BEGINS
% ============================================
\begin{document}

% ============================================
% TITLE SLIDE (Plain, White, Centered)
% ============================================
\begin{frame}[plain]
    \begin{center}
        \vspace{0.8cm}
        {\Large \bfseries \ProjectTitle \par}
        \vspace{1.5em}
        {\normalsize \bfseries Team Members:\par}
        \vspace{0.3em}
        {\small \studentone{} (\rollone)\par}
        {\small \studenttwo{} (\rolltwo)\par}
        {\small \studentthree{} (\rollthree)\par}
        \vfill
        {\normalsize \DepartmentName\par}
        {\normalsize \InstitutionName\par}
        \vspace{0.3em}
        {\normalsize \PresentationDate\par}
    \end{center}
\end{frame}

% ============================================
% PRESENTATION OUTLINE
% ============================================
\begin{frame}{Presentation Outline}
    \begin{itemize}
        \item Origins and Philosophy of Lean
        \item Traditional Problems and Solution
        \item The Seven Process Wastes
        \item The 7 Core Lean Principles
        \item Tools and Practical Comparison
        \item Advantages, Disadvantages, and Case Study
    \end{itemize}
\end{frame}

% ============================================
% SECTION: ORIGINS AND PHILOSOPHY (Member 1)
% ============================================

% Page 2
\begin{frame}{What is Lean Software Development?}
    Lean Software Development is a software development approach that focuses on delivering maximum value to customers while minimizing waste.
    \vspace{0.8em}
    
    \textbf{Key Points}
    \begin{itemize}
        \item Inspired by Lean Manufacturing and the Toyota Production System
        \item Adapted for software development by Mary and Tom Poppendieck
        \item Emphasizes efficiency, quality, and customer satisfaction
        \item Supports Agile and Continuous Delivery practice
    \end{itemize}
\end{frame}

% Page 3
\begin{frame}{Why Lean Was Introduced}
    \textbf{Problems with Traditional Development}
    \vspace{0.5em}
    
    \begin{itemize}
        \item Traditional software development required extensive documentation before coding began.
        \item A large amount of time and effort was spent creating documents.
        \item Documents often became outdated as project requirements changed.
        \item This reduced efficiency and slowed down software delivery.
        \item Feedback was often received late in the development process.
        \item Late changes increased project cost, effort, and delays.
    \end{itemize}
\end{frame}

% Page 4
\begin{frame}{Lean Solution}
    \begin{itemize}
        \item Lean focuses on delivering value to customers as quickly as possible.
        \item Working software is released in smaller and more frequent increments.
        \item Continuous feedback helps teams improve and adapt to changes.
        \item Waste is minimized, improving overall efficiency and productivity.
    \end{itemize}
\end{frame}

% Page 5
\begin{frame}{The Seven Wastes in Software Development}
    \textbf{Seven Types of Waste}
    \vspace{0.4em}
    
    \begin{enumerate}
        \item \textbf{Partially Done Work} - Incomplete tasks that are not yet useful
        \item \textbf{Extra Features} - Features that customers do not need
        \item \textbf{Relearning} - Loss of knowledge requiring repeated effort
        \item \textbf{Task Switching} - Reduced productivity from multitasking
        \item \textbf{Waiting} - Delays between development activities
        \item \textbf{Handoffs} - Information loss during communication
        \item \textbf{Defects} - Errors that require correction and rework
    \end{enumerate}
\end{frame}

% ============================================
% SECTION: THE 7 PRINCIPLES (Member 2)
% ============================================

% Page 7
\begin{frame}[plain]
    \begin{center}
        \vspace{2cm}
        {\Large \bfseries The 7 principles of lean software development}
    \end{center}
\end{frame}

% Page 8
\begin{frame}{Eliminate Waste}
    \begin{center}
        \textit{If it doesn't add value for the user --- remove it.}
    \end{center}
    \vspace{0.8em}
    
    \begin{itemize}
        \item Remove unused features, duplicate steps, and unnecessary waiting
        \item Avoid writing code that nobody asked for
        \item Less clutter means faster delivery and a cleaner codebase
        \item Waste includes: delays, defects, extra processes, and over-engineering
    \end{itemize}
\end{frame}

% Page 9
\begin{frame}{Build Quality In}
    \begin{center}
        \textit{Catch problems early --- prevention is cheaper than repair.}
    \end{center}
    \vspace{0.8em}
    
    \begin{itemize}
        \item Test code as it is written, not after the entire product is done
        \item Use automated test suites so every change is verified instantly
        \item Code reviews and pair programming catch issues before they ship
        \item A bug fixed in development costs $10\times$ less than one found in production
    \end{itemize}
\end{frame}

% Page 10
\begin{frame}{Create Knowledge}
    \begin{center}
        \textit{Learn continuously and share what you know.}
    \end{center}
    \vspace{0.8em}
    
    \begin{itemize}
        \item Document decisions and the reasons behind them
        \item Hold short retrospectives after each sprint or milestone
        \item Maintain a shared wiki so lessons reach the whole team
        \item Don't repeat the same mistakes i.e. knowledge compounds over time
    \end{itemize}
\end{frame}

% Page 11
\begin{frame}{Defer Commitment}
    \begin{center}
        \textit{Delay big decisions until you have enough information.}
    \end{center}
    \vspace{0.8em}
    
    \begin{itemize}
        \item Keep options open as long as it's practical to do so
        \item Avoid expensive, hard-to-reverse choices made on early assumptions
        \item Gather data and user feedback before locking in an architecture
        \item Decide at 'the last responsible moment' --- not the last possible moment
    \end{itemize}
\end{frame}

% Page 12
\begin{frame}{Deliver Fast}
    \begin{center}
        \textit{Release early, release often - get feedback sooner.}
    \end{center}
    \vspace{0.8em}
    
    \begin{itemize}
        \item Break large projects into short, shippable work cycles
        \item Deliver small updates rather than one massive release
        \item Fast delivery surfaces problems before they snowball
        \item Speed reduces risk i.e. smaller batches are easier to test and roll back
    \end{itemize}
\end{frame}

% Page 13
\begin{frame}{Respect People}
    \begin{center}
        \textit{Trust your team and give them real ownership.}
    \end{center}
    \vspace{0.8em}
    
    \begin{itemize}
        \item Let developers make decisions about their own work
        \item Listen to feedback and act on what the team surfaces
        \item Protect work-life balance to sustain long-term performance
        \item Psychological safety drives creativity and honest communication
    \end{itemize}
\end{frame}

% Page 14
\begin{frame}{Optimize the Whole}
    \begin{center}
        \textit{Look at the entire pipeline, not just one step.}
    \end{center}
    \vspace{0.8em}
    
    \begin{itemize}
        \item Speeding up one stage can create bottlenecks elsewhere
        \item Map the full journey from idea to production deployment
        \item Remove handoff delays between teams (dev $\to$ QA $\to$ ops)
        \item Global optimisation beats local improvements every time
    \end{itemize}
\end{frame}

% ============================================
% SECTION: PRACTICES AND TOOLS (Member 3)
% ============================================

% Page 16
\begin{frame}{Tools \& Techniques in Lean}
    \textbf{Value Stream Mapping (VSM)}
    \begin{itemize}
        \item Diagrams every step from requirement to delivery
        \item Highlights delays, waste, and unnecessary handoffs
    \end{itemize}
    \vspace{0.3em}
    
    \textbf{Kanban Board}
    \begin{itemize}
        \item Visualizes tasks: To Do, In Progress, Done
        \item Limits work-in-progress to prevent overload
    \end{itemize}
    \vspace{0.3em}
    
    \textbf{Pull-Based Scheduling}
    \begin{itemize}
        \item Work begins only when team has capacity
        \item Prevents bottlenecks from piling up downstream
    \end{itemize}
    \vspace{0.3em}
    
    \textbf{Continuous Integration}
    \begin{itemize}
        \item Code tested automatically after every change
        \item Defects caught immediately, not at release
    \end{itemize}
\end{frame}

% Page 17
\begin{frame}{Lean vs Agile Development}
    Lean is a philosophy; Agile is a methodology
    \vspace{0.3em}
    
    Both reject heavy upfront planning and documentation
    \vspace{0.3em}
    
    Agile Manifesto (2001) was shaped by Lean thinking
    \vspace{0.3em}
    
    Scrum implements Lean through sprints and roles
    \begin{itemize}
        \item Roles: Product Owner, Scrum Master, Dev Team
        \item Fixed iterations: typically 1--4 week sprints
    \end{itemize}
    \vspace{0.3em}
    Lean has no fixed iterations or prescribed roles
    \vspace{0.3em}
    
    Key shared goal: deliver working software faster
\end{frame}

% Page 18
\begin{frame}{Advantages}
    \begin{itemize}
        \item Reduces time from requirement to working software
        \item Defects caught early cost significantly less to fix
        \item Removes steps that add no customer value
        \item Frequent releases allow continuous user feedback
        \item Developers given autonomy produce higher quality work
        \item Cleaner code reduces long-term maintenance cost
    \end{itemize}
\end{frame}

% Page 19
\begin{frame}{Disadvantages}
    \begin{itemize}
        \item Requires full organizational commitment to work
        \item Difficult without experienced and disciplined team members
        \item Reduced documentation risks long-term maintainability
        \begin{itemize}
            \item System knowledge becomes dependent on individuals
            \item New members struggle without proper documentation
        \end{itemize}
        \item Over-reliance on key members increases project risk
        \item Progress and improvement difficult to measure objectively
    \end{itemize}
\end{frame}

% Page 20
\begin{frame}{Case Study: Toyota Production System}
    \begin{itemize}
        \item Large batch production hid defects until very late
        \item Value stream mapping revealed primary delay points
        \begin{itemize}
            \item Handoffs between stations were main bottleneck
        \end{itemize}
        \item Kanban introduced to pull work by actual capacity
        \item Defects surfaced and resolved immediately at each stage
        \item Software teams later applied same principles to code delivery
    \end{itemize}
\end{frame}

% Page 21
\begin{frame}{Conclusion}
    \begin{itemize}
        \item Lean targets root causes of software project failure
        \item Waste elimination improves both speed and quality
        \item Agile and Lean share the same foundational goals
        \item Documentation remains critical despite Lean's speed focus
        \item Team capability determines how well Lean is applied
        \item Lean thinking applies across any development context
    \end{itemize}
\end{frame}

\begin{frame}[plain]
    \begin{center}
        \vspace{1.5cm}
        {\huge \bfseries Thank You} \\[1em]
        {\large Questions \& Answers}
    \end{center}
\end{frame}

\end{document}
